\section{Sécurité, intégrité et accessibilité}

\subsection{Modèle de menace}

Le risque principal n’est pas l’interception d’une requête de traduction :
KORA n’en émet aucune. Le risque est une modification silencieuse du sens ou
d’une valeur canonique pendant le rendu. La conception traite donc la
traduction comme une transformation non fiable qui doit prouver qu’elle a
conservé les données.

\begin{tabularx}{\linewidth}{@{}|L{40mm}|Y|Y|@{}}
\hline
\textbf{Risque} & \textbf{Impact possible} & \textbf{Mesure KORA}\\
\hline
Montant ou devise altéré &
Décision financière erronée &
Masquage, restauration exacte, comptage des occurrences, rejet.\\
\hline
Bénéficiaire traduit &
Confusion sur le destinataire &
Noms et valeurs saisies ajoutés au jeu protégé en mémoire d’onglet.\\
\hline
Action inversée &
Confirmer au lieu d’annuler, créditer au lieu de débiter &
Actions critiques en mémoire exacte; négation inconnue = abstention.\\
\hline
Alerte affaiblie &
Utilisateur insuffisamment averti &
Messages complets, seuil de couverture critique à 100\,\%.\\
\hline
Fuite de contenu &
Exposition de données personnelles &
Pas de transport réseau, télémétrie ou stockage persistant de page.\\
\hline
Régression d’interface &
Clavier ou formulaire inutilisable &
Modification des nœuds textuels, pas des composants ou de leur ordre.\\
\hline
\end{tabularx}

\subsection{Invariant d’intégrité}

Soit \(P(s)\) le multiensemble des valeurs protégées extraites d’un segment
source \(s\), et \(t\) une traduction candidate. KORA accepte \(t\) seulement
si :

\[
  \forall v \in P(s),\quad
  \operatorname{occ}(v,s)=\operatorname{occ}(v,t)
  \quad\text{et}\quad
  \operatorname{bytes}(v_s)=\operatorname{bytes}(v_t).
\]

En pratique, les valeurs sont remplacées temporairement par des jetons
internes, par exemple \code{<<KORA\_0>>}, puis remises caractère pour caractère.
Le contrôle final protège aussi les répétitions : une valeur ne doit ni
disparaître, ni être dupliquée.

\begin{keybox}[Exemple de transformation sûre]
\textbf{Source :} Votre transfert de \texttt{25 000 FCFA} est en attente.\\
\textbf{Interne :} Votre transfert de \texttt{<<KORA\_0>>} est en attente.\\
\textbf{Candidat :} Transfeer maa \texttt{<<KORA\_0>>} ɗon e fadde.\\
\textbf{Rendu :} Transfeer maa \texttt{25 000 FCFA} ɗon e fadde.
\end{keybox}

\subsection{Données protégées}

\begin{multicols}{2}
\begin{itemize}
  \item montants, signes et séparateurs;
  \item monnaies et pourcentages;
  \item comptes et numéros de téléphone;
  \item références et codes d’erreur;
  \item dates, heures, URL et courriels;
  \item noms complets et valeurs utilisateur.
\end{itemize}

\columnbreak
\begin{itemize}
  \item \code{input.value} et \code{textarea.value};
  \item \code{select.value} et valeur technique d’une option;
  \item \code{checked}, \code{selected}, \code{name}, \code{id};
  \item \code{type}, \code{min}, \code{step};
  \item contenu \code{contenteditable};
  \item identifiants opaques comme \code{DSAFDF}.
\end{itemize}
\end{multicols}

Dans un \code{select}, seul le libellé visible d’une option peut changer.
Sa valeur technique reste intacte. Le texte saisi
\texttt{Payment for services} reste lui aussi intact, même si le label du
champ est traduit.

\subsection{Cycle de vie des données}

\begin{figure}[H]
\centering
\begin{tikzpicture}[
  node distance=6mm and 12mm,
  b/.style={draw=Line,fill=Paper,rounded corners=1.5mm,
    text width=31mm,minimum height=11mm,align=center,font=\small},
  ok/.style={b,draw=Teal,fill=Mint},
  no/.style={b,draw=Gold,fill=Warm},
  arrow/.style={-{Latex[length=2mm]},thick,color=TealDark}
]
\node[b] (read) {Lecture temporaire du segment visible};
\node[b,right=of read] (protect) {Extraction des valeurs et masquage};
\node[b,right=of protect] (translate) {Traitement local en mémoire};
\node[ok,right=of translate] (render) {Rendu du texte ou original};
\node[no,below=9mm of protect] (store) {Jamais écrit dans le stockage};
\node[no,below=9mm of translate] (network) {Jamais envoyé sur le réseau};
\draw[arrow] (read)--(protect);
\draw[arrow] (protect)--(translate);
\draw[arrow] (translate)--(render);
\draw[-{Latex[length=2mm]},thick,color=Risk] (protect)--(store);
\draw[-{Latex[length=2mm]},thick,color=Risk] (translate)--(network);
\end{tikzpicture}
\caption{Le contenu financier reste dans la mémoire volatile de l’onglet.}
\end{figure}

\code{chrome.storage.local} conserve uniquement la préférence d’activation et
la langue d’affichage. Le jeu de valeurs utilisateur, utilisé pour reconnaître
le même nom ou la même note au récapitulatif, existe uniquement en mémoire du
content script et disparaît avec l’onglet.

\subsection{Permissions et indépendance}

\begin{lstlisting}[language={},caption={Permissions déclarées dans manifest.json}]
"permissions": ["storage", "activeTab", "scripting"]
\end{lstlisting}

\begin{itemize}
  \item \code{storage} conserve les préférences non sensibles.
  \item \code{activeTab} relie la popup à l’onglet explicitement utilisé.
  \item \code{scripting} charge le moteur local si l’onglet ne possède pas
        encore le protocole courant.
\end{itemize}

Le content script est déclaré pour les pages ordinaires HTTP, HTTPS et fichiers
locaux. Cette portée démontre l’indépendance par rapport à Safinum, mais KORA
ne contourne ni politique de contenu, ni iframe interdite, ni page interne
Chrome.

\subsection{Accessibilité et expérience utilisateur}

\begin{tabularx}{\linewidth}{@{}|L{43mm}|Y|@{}}
\hline
\textbf{Principe} & \textbf{Application}\\
\hline
Navigation clavier &
Aucun gestionnaire global ne bloque \code{Ctrl+A}; tabulation et commandes du
site sont conservées.\\
\hline
Focus &
La popup possède des styles de focus visibles et se ferme après action pour
rendre le focus à la page.\\
\hline
Sémantique &
Les éléments ne sont ni clonés ni remplacés; rôles, labels et ordre DOM
restent ceux du site.\\
\hline
Lecteurs d’écran &
\code{aria-label}, \code{title} et \code{placeholder} peuvent être traduits;
les sorties portent \code{lang="fub"}.\\
\hline
Lisibilité &
Popup sombre, contrastes élevés, contrôles dimensionnés et états explicites.\\
\hline
Réversibilité &
Chaque nœud conserve sa forme source; l’utilisateur peut restaurer la page
instantanément.\\
\hline
\end{tabularx}

\subsection{Cadre de démonstration}

Toutes les opérations sont simulées. Les comptes, noms, numéros, montants et
références du parcours servent uniquement à vérifier l’intégrité; KORA ne
déclenche aucun transfert et ne doit recevoir aucune donnée bancaire réelle.
La démonstration respecte la page fournie : elle charge l’extension comme une
extension utilisateur, sans modifier le code source ni contourner les étapes.

\begin{warningbox}[Limite réglementaire]
KORA est un prototype d’accessibilité, pas un dispositif bancaire certifié,
un outil d’autorisation de transaction ou un conseil financier. Un déploiement
réel exigerait une revue de sécurité indépendante, une double validation
linguistique, des essais utilisateurs et la gouvernance versionnée du
glossaire.
\end{warningbox}
